\section{SpartinaShift-Silver v0.1: frozen-protocol conventional baselines}\label{sec:silverbench}

This section reports the first model benchmark executed on the real
Issue~\#19 V2 pixels: 20 fixed W10 cells, 184 landed products, actual
in-cell valid masks. The benchmark is SILVER-only throughout (GOLD\,=\,0),
and every metric is a statement about agreement with a named label
source, never about mapping accuracy. The protocol was frozen before any
training run (\texttt{SPARTINASHIFT\_SILVER\_V0\_1\_PROTOCOL}, frozen
2026-10-12, SHA-256 \texttt{24037877\ldots}); its checksum, the split
and support-matrix hashes, the git commit, seed, full configuration,
checkpoint hash and software environment are recorded in every run
manifest. Rights follow the frozen policy: GEODATA layers are
\texttt{SILVER\_TRAINING\_ELIGIBLE\_INTERNAL\_RESEARCH\_ONLY}, CMSA is
\texttt{SILVER\_TRAINING\_ELIGIBLE\_NONCOMMERCIAL}, and CM-SSM is
\texttt{EXTERNAL\_REFERENCE\_ONLY} under a source-license conflict---it
is never trained or tuned on. The finding is an honest negative. On
spatially disjoint held-out cells every conventional baseline performs
at chance after a threshold frozen on validation data, in both
evaluation years, with or without SAR. Diagnostic controls show that the
within-scene spectral signal is real and the pipeline recovers it; what
collapses is transfer, and the collapse is localised to measured support
failures rather than to model capacity.

\subsection{Frozen design: cell-disjoint split, 30\,m support, locked selection}
The W10 Albers 10\,km cell is the only split unit. Patches, dates,
modalities and years of one cell inherit its role, cross-role cell
footprints are separated by a 250\,m buffer (the measured minimum
cross-role gap is 80\,km), and the split---11 train, 3 validation, 4
test cells, with two named 2020 witness cells inside the test set---was
drawn once and never reseeded. Training uses CMSA 2020 labels on
Sentinel-2 B2/B3/B4/B8; tests are opened on CMSA 2020 (reference year)
and CMSA 2021 (cross-year). Labels follow the fractional rule: a pixel
is binary-positive at mapped fraction $\geq$0.90, MIXED pixels (0 to
0.90, product-specific lower bound) are excluded from both the binary
loss and the primary metrics and are described separately, and unknown
or EO-invalid pixels are ignored. SILVER background is never treated as
a verified ecological negative.

All observations are analysed on the 30\,m label grid: native 10\,m
Sentinel-1/2 data are averaged down with a sub-pixel validity rule, and
no label pixel is ever upsampled. Patches are 128\,px (3.84\,km on the
ground), 128 per training cell per epoch with a 0.5 positive-patch
fraction, with six deterministic geometric augmentations. Per-band
p005/p995 clip bounds and mean/std come from valid pixels of train cells
only, separately per arm. Hyperparameters are fixed by the protocol
with no search: a pixel Random Forest (300 trees, 20{,}000-pixel cap per
train cell, \texttt{sqrt} features, minimum leaf 2, EO bands plus NDVI;
the SAR arm adds VV, VH and the VV$-$VH dB difference), U-Net and
DeepLabV3+ with ResNet18 ImageNet encoders (40 epochs), and
SegFormer-B0 with a replaced segmentation head (30 epochs); networks use
0.5 BCE plus 0.5 soft-Dice loss with positive weight capped at 20,
AdamW with cosine schedule and three warm-up epochs, and mixed
precision. Each run uses seeds 17, 42 and 2026. One decision threshold
per run is selected on pooled validation IoU over a fixed 0.05--0.95
grid, with patience-10 early stopping on the same quantity; the test
cells are opened only after the checkpoint and threshold are frozen.
Uncertainty is reported as per-cell values on independent cells, plus a
descriptive 1\,000\,m spatial-block bootstrap (1\,000 replicates, blocks
never cross cell boundaries); with only two positive test cells, no
inferential cell-level claim is made.

The support itself is the first result (Table~\ref{tab:silversupport}).
Only five of the 20 cells contain any CMSA-2020 pure-positive pixels:
two train, one validation and two test cells. One train cell has
virtually no within-scene NDVI separation, and one test cell has zero
EO-valid positive pixels in 2020. The executed campaign comprises 24
training runs (tasks C2 and C5, four models by three seeds) and 12
frozen-checkpoint agreement evaluations (task C4); tasks C1 and C3
executed zero runs by pre-registered support rule.

% Source: benchmarks/spartinashift/silver_v0_1/results/C2/failure_analysis_C2.json
% and SPARTINASHIFT_SILVER_V0_1_PROTOCOL.json (frozen, checksum
% 24037877...). CMSA 2020 is the C2 training label; GOLD = 0.
\begin{table}[t]
\centering
\small
\caption{SpartinaShift-Silver v0.1 (SPLIT-S2) binary-positive support on
the 20-cell Issue~\#19 V2 panel at 30\,m. Only five cells contain
CMSA-2020 pure-positive pixels (fraction $\geq$0.90): two train, one
validation, two test. EO-valid positives intersect the actual Sentinel-2
valid mask inside each cell; the NDVI gap is positive-minus-background
mean NDVI on 2020 valid pixels. All other 15 cells are source-background
only, and SILVER background is never a verified ecological negative.}
\label{tab:silversupport}
\begin{tabular}{llrrrr}
\toprule
Role & Cell (province) & Pure+ px & EO-valid+ & EO-valid \% & NDVI gap \\
\midrule
train & R00311 (Zhejiang) & 3{,}981 & 784 & 19.7 & $-$0.005 \\
train & R00349 (Jiangsu) & 5{,}849 & 3{,}257 & 55.7 & +0.532 \\
val   & R00236 (Guangdong) & 149 & 135 & 90.6 & +0.104 \\
test  & R00334 (Zhejiang) & 14 & 0 & 0.0 & n/a \\
test  & R00413 (Shandong) & 1{,}335 & 1{,}322 & 99.0 & +0.486 \\
\midrule
\multicolumn{6}{l}{Other 15 cells (9 train, 2 val, 2 test): 0 pure-positive pixels} \\
\bottomrule
\end{tabular}
\end{table}


\subsection{Task C2: cross-year Sentinel-2 baselines at chance}
Table~\ref{tab:silverc2} reports pooled held-out IoU and F1. No model
reaches a usable operating point. Pooled test IoU is at most 0.003 in
2020 and 0.020 in 2021, and the latter value is a single DeepLabV3+
seed (mean 0.020, SD 0.028); U-Net scores exactly zero in both years on
every seed. Every spatial-block 95\% interval includes zero.
Validation is already weak---the best mean validation IoU is 0.038
(U-Net) on a single positive validation cell containing 149 positive
pixels---so the frozen thresholds are selected on almost no signal.
Rank information nonetheless survives where positives exist: on cell
R00413 in 2021 the per-cell AUPRC is 0.297 for one U-Net run even
though thresholded IoU at its frozen validation threshold is 0.000,
which locates the failure partly in threshold transfer rather than in
wholly uninformative scores. The networks also issue false alarms on
the two background-only test cells (e.g.\ 447\,ha of predicted positive
area there for one U-Net seed in 2020). Seed spread is large relative
to every mean, and MIXED-pixel probability--fraction rank correlations
stay weak (0.16 in 2020 and 0.24 in 2021 for the same run), reported
descriptively only.

% Source: benchmarks/spartinashift/silver_v0_1/results/C2/*_s2_seed*.json
% and aggregate_C2_s2.json (frozen protocol checksum 24037877...;
% split SPLIT-S2; git 66549b0). 12 runs = 4 models x seeds 17/42/2026.
\begin{table}[t]
\centering
\small
\caption{Task C2: CMSA-2020 $\rightarrow$ 2020/2021 cross-year
Sentinel-2 baselines, spatially disjoint W10 cells (11/3/4; only two
positive test cells). Values are pooled IoU/F1 over the two positive
held-out cells, mean$\pm$SD over the three seeds. The bracketed interval
is the union across seeds of the 1\,000\,m spatial-block bootstrap 95\%
intervals (1\,000 replicates; blocks never cross cells; 15--411 blocks
per seed). VAL IoU is pooled over the single positive validation cell
(R00236, 149 positive pixels); the decision threshold is chosen on VAL
only, once per run. Every TEST 95\% interval includes 0.}
\label{tab:silverc2}
\begin{tabular}{lrrrlrl}
\toprule
& \multicolumn{1}{c}{VAL 2020} & \multicolumn{3}{c}{TEST 2020} & \multicolumn{2}{c}{TEST 2021} \\
\cmidrule(lr){2-2}\cmidrule(lr){3-5}\cmidrule(lr){6-7}
Model & IoU & IoU & block CI95 & F1 & IoU & F1 \\
\midrule
RF pixel      & 0.0011\,$\pm$\,0.0001 & 0.0013\,$\pm$\,0.0001 & [0.0003, 0.0027] & 0.0026 & 0.0001 & 0.0002 \\
U-Net         & 0.0378\,$\pm$\,0.0131 & 0.0000 & [0, 0] & 0.0000 & 0.0000 & 0.0000 \\
DeepLabV3+    & 0.0250\,$\pm$\,0.0180 & 0.0004\,$\pm$\,0.0005 & [0, 0.0034] & 0.0008 & 0.0202\,$\pm$\,0.0275 & 0.0382 \\
SegFormer-B0  & 0.0330\,$\pm$\,0.0110 & 0.0033\,$\pm$\,0.0036 & [0, 0.0222] & 0.0065 & 0.0004\,$\pm$\,0.0006 & 0.0008 \\
\bottomrule
\end{tabular}
\end{table}


\subsection{Why the models fail: same-scene signal, no transferable rule}
Four diagnostic controls, computed once and versioned with the
benchmark, separate ``no signal'' from ``no transfer''. None uses test
selection.

First, the same-scene ceiling is high. A Random Forest trained and
thresholded within one scene reaches IoU 0.388, 0.478, 0.430 and 0.630
on the four positive cells with usable overlap (R00236, R00311, R00349,
R00413); R00334 is not estimable because it has zero EO-valid positive
pixels. Where a scene contains separated Spartina cover, the features,
labels and evaluation machinery recover it.

Second, cross-scene transfer collapses even between the two strongly
separable scenes. Training on R00349 and evaluating R00413 gives IoU
0.035, and the reverse gives 0.027; the two scenes differ in absolute
reflectance even though each is internally separable, so a decision
rule learned on one does not carry to the other.

Third, temporal transfer inside the same scene is weak: a 2020 model
applied to the same cell in 2021 reaches IoU 0.176 on R00349 and 0.024
on R00413. Fourth, the acquisition calendar is single-date: each
cell-year contains one autumn Sentinel-2 event, and on R00413 the two
dates are 2020-09-18 and 2021-11-12---day-of-year 262 versus 316, 420
calendar days apart. The nominal ``cross-year'' test therefore embeds a
large phenology and acquisition shift that a single-date training set
cannot span.

The support pathologies compound this. The only positive training cell
with wide coverage (R00311) has just 19.7\,\% of its pure-positive
pixels under valid EO and an NDVI gap of $-$0.005; the other strong
train cell (R00349) supplies most of the usable signal but does not
transfer; the validation set contains one positive cell with 149
pixels; and the test cell R00334 contributes 14 positive pixels in 2020,
none of them EO-valid, and 8 in 2021. A pipeline explanation for the
near-zero scores is thus not supported by the controls; a support
explanation is.

\subsection{Tasks C1 and C3: pre-registered non-runs}
Two planned tasks were refused before training by the frozen protocol,
and their records contain no metrics beyond
\texttt{NOT\_RUN\_INSUFFICIENT\_SUPPORT}. C1 (GEODATA cross-year on
Landsat~8, 2015$\rightarrow$2020) requires at least five positive cells
per year for a spatial binary split; GEODATA is sparse inside the
panel, with three pure-positive cells in 2015 and three in 2020, only
two of them positive in both years, even though 14 cells are EO-eligible
in both years. C3 (cross-region) fails its frozen minimums in all seven
provinces present---at least four held-out cells with two positive, and
eight remaining cells with three positive outside---including its best
case, Zhejiang, which holds two positive cells but only two cells in
total. No regional split was manufactured by relaxing these rules.

\subsection{Task C4: cross-source agreement at 2020, without tuning}
C4 reuses each frozen C2 checkpoint, its training-only normalization
and its validation-selected threshold verbatim: no retraining, no
threshold adjustment, and CM-SSM is never used for model selection in
any form. Agreement is then measured against GEODATA 2020 on the two
positive test cells and against CM-SSM 2020 on those cells plus the two
named witness cells (Table~\ref{tab:silverc4}). Agreement is negligible
throughout, and the label sources disagree with each other at cell
scale. On R00413 CMSA maps 1{,}335 pure-positive pixels and GEODATA
1{,}542, while CM-SSM maps zero; on R00334 GEODATA maps 40 and CM-SSM
136. The witness background cell R00445 is zero in both sources. With
no GOLD labels this disagreement cannot be adjudicated---near-zero
agreement of a CMSA-trained model with CM-SSM is neither a model error
nor a CM-SSM error---and the result is reported strictly as
external-reference agreement under the CM-SSM license restriction. The
Random Forest GEODATA interval narrowly excludes zero
([0.0022, 0.0150]) at an IoU magnitude of 0.008; nothing else excludes
zero.

% Source: benchmarks/spartinashift/silver_v0_1/results/C4/*.json and
% aggregate_C4.json. Frozen C2 checkpoints and VAL-only thresholds are
% reused; no retraining and no tuning on either 2020 label source.
% Agreement metrics, NOT accuracy: GOLD = 0. CM-SSM is
% EXTERNAL_REFERENCE_ONLY (source license conflict).
\begin{table}[t]
\centering
\small
\caption{Task C4: agreement of frozen CMSA-trained C2 models with two
independent 2020 SILVER label sources---GEODATA 2020 on the two positive
test cells (cross-source agreement), and CM-SSM 2020 on those cells plus
the two named witness cells (external-reference agreement; CM-SSM was
never trained or tuned). Mean$\pm$SD over seeds 17/42/2026; bracketed
intervals are the union of per-seed 1\,000\,m spatial-block 95\% CIs.
Agreement is not accuracy: neither source is verified truth.}
\label{tab:silverc4}
\begin{tabular}{lrrlrl}
\toprule
& \multicolumn{3}{c}{GEODATA 2020} & \multicolumn{2}{c}{CM-SSM 2020} \\
\cmidrule(lr){2-4}\cmidrule(lr){5-6}
Model & IoU & block CI95 & F1 & IoU & F1 \\
\midrule
RF pixel      & 0.0078\,$\pm$\,0.0005 & [0.0022, 0.0150] & 0.0155 & 0.0008 & 0.0016 \\
U-Net         & 0.0000 & [0, 0] & 0.0000 & 0.0007\,$\pm$\,0.0005 & 0.0013 \\
DeepLabV3+    & 0.0000 & [0, 0] & 0.0000 & 0.0000 & 0.0000 \\
SegFormer-B0  & 0.0027\,$\pm$\,0.0039 & [0, 0.0244] & 0.0054 & 0.0019\,$\pm$\,0.0026 & 0.0037 \\
\bottomrule
\end{tabular}
\end{table}


\subsection{Task C5: adding Sentinel-1 does not rescue the support}
The modality ablation repeats the ladder with six channels
(S2 B2/B3/B4/B8 plus S1-A VV/VH) on the 17 cells with paired products
(one Fujian cell is excluded for lack of S1-A; roles are 10/3/4), with
identical held-out cells, frozen patch geometry, seeds and budget, and
per-arm train-only normalization. Table~\ref{tab:silverc5} shows the
comparison. Movement is within seed noise and in both directions: U-Net
rises from exactly zero to 0.016 mean IoU in 2020 (SD 0.009) while
DeepLabV3+ falls from 0.020 to 0.003 in 2021; no arm produces a stable
positive result. The pairing itself is season-matched, not
date-matched: all 34 cell-year pairs fall inside the 30-day window but
the median gap is 22.5 days (range 1--30), only five pairs are within
seven days and none is simultaneous, and the five-pair $\leq$7-day
subset is pre-registered as insufficient for its own inference. The
extra channels also damage stability---mean expected calibration error
for DeepLabV3+ on 2020 test pixels rises from 0.009 to 0.111 and its
mean predicted positive area on background cells rises from 333 to
9{,}641\,ha---while U-Net validation IoU rises to 0.133 (SD 0.091)
without any matching test gain, the same single-cell validation
overfitting seen in C2.

% Source: benchmarks/spartinashift/silver_v0_1/results/C5/*.json,
% aggregate_C5_s2_s1.json and date_gap_ledger.json. 12 new runs on the
% 17 S2/S1-A paired cells (R00279 excluded, no S1-A product).
\begin{table}[t]
\centering
\small
\caption{Task C5 modality ablation: pooled held-out IoU, optical-only
(S2, 4\,ch, frozen C2 runs) versus optical+SAR (S2+S1-A, 6\,ch, 12 new
runs), identical held-out cells, frozen patch geometry, seeds and
training budget; per-arm train-only normalization. Mean$\pm$SD over
seeds 17/42/2026. The 34 S2--S1-A cell-year pairs are all within the
30-day season window (median gap 22.5\,d, range 1--30\,d; five within
7\,d; zero simultaneous), so the comparison is season-matched rather
than date-matched. Adding SAR moves no model off chance on the held-out
cells.}
\label{tab:silverc5}
\begin{tabular}{lcccc}
\toprule
& \multicolumn{2}{c}{TEST 2020 IoU} & \multicolumn{2}{c}{TEST 2021 IoU} \\
\cmidrule(lr){2-3}\cmidrule(lr){4-5}
Model & S2 & S2+S1-A & S2 & S2+S1-A \\
\midrule
RF pixel      & 0.0013 & 0.0059\,$\pm$\,0.0001 & 0.0001 & 0.0001 \\
U-Net         & 0.0000 & 0.0158\,$\pm$\,0.0085 & 0.0000 & 0.0071\,$\pm$\,0.0065 \\
DeepLabV3+    & 0.0004 & 0.0028\,$\pm$\,0.0018 & 0.0202 & 0.0025\,$\pm$\,0.0018 \\
SegFormer-B0  & 0.0033 & 0.0036\,$\pm$\,0.0032 & 0.0004 & 0.0070\,$\pm$\,0.0020 \\
\bottomrule
\end{tabular}
\end{table}


\subsection{What the v0.1 benchmark establishes---and what it does not}
Five statements are supported by executed runs. The frozen protocol
runs end to end with complete provenance and leakage-safe cell
splits; conventional pixel and segmentation baselines are at chance on
spatially disjoint cells at this support level; within-scene spectral
signal demonstrably exists, so the null result is a transfer failure;
the transfer failure is explained by measured mechanisms---sparse and
partly invalid positive support, a zero-separation training cell,
single-date acquisitions with large inter-year phenology offsets, and
substantial disagreement among independent SILVER sources; and adding
Sentinel-1 under season-matched pairing does not repair any of them.
Equally explicitly, the benchmark establishes no accuracy claim
(GOLD\,=\,0), nothing about the intrinsic cross-region
transferability of \emph{Spartina alterniflora} spectra, and nothing
about whether higher-capacity models would succeed---the experiment was
not run, and the diagnosed bottlenecks (zero valid positives, one
separable cross-scene pair, 22-day median optical--SAR gaps) are not
capacity bottlenecks. Twenty cells and two positive test cells are a
pilot support for freezing protocol and instrumentation, not benchmark
v1. The defensible next gate is support repair---more positive cells,
multi-date seasonal composites rather than one event per year, denser
near-simultaneous optical--SAR pairing, and GOLD-verified labels---before
model capacity is increased; the corresponding Stage-1 readiness token
is recorded with the benchmark registry and issue report rather than
inferred from these numbers.

\begin{figure}[t]\centering
\figplaceholder{Figure 8}{SpartinaShift-Silver v0.1 evidence panel
(design only): (a)~SPLIT-S2 cell roles on the 20-cell panel with the
five CMSA-positive cells; (b)~same-scene IoU ceiling (0.39--0.63)
versus cross-scene (0.035/0.027) and temporal transfer (0.176/0.024);
(c)~held-out IoU of the four models, both years, with spatial-block
95\% intervals all including zero; (d)~2020 per-cell positive counts
for CMSA, GEODATA and CM-SSM showing source disagreement. All panels
annotated GOLD\,=\,0 and ``agreement, not accuracy''.}
\label{fig:silverbench}
\end{figure}
